We report Recall@1 (R@1), Recall@3 (R@3) and Normalized Discounted Cumulative Gain (nDCG@3) metrics on the test set for each book. BeIR~\cite{thakur2021beir} is used for computing the metrics.

% \eat{
% Our goal is to show the efficacy of {\sys} in identifying the correct subsection ({\toc} leaf node) of the document for the given query. We use proposed {\ben} benchmark dataset (described in Section ~\ref{sec:benchmark}) to compare performance of {\sys} with other strong baselines like BM25, DSI etc. We use the train set of the {\ben} dataset for supervised finetuning, use recall@1 score on the dev set for model checkpoint selection and report the final performance on the test set. The subsequent sub-sections will outline the specifics of the baseline models, prompts, supervised fine-tuning procedures, and evaluation metrics.
% }

% For baselines such as DPR or DSI When we have multiple trained checkpoints, we chose the model which gives us the best Recall@1 numbers on the development set and use it to report the final numbers on the test set. This setup is well aligned with how we evaluate {\sys} as well. \eat{In order to build (or train) an index, we only use the train questions. For some retrievers we have multiple options to select from -- we chose the model which gives us the best Recall@1 numbers on the development set. The best model is then used to report the final numbers on the test set.} Next, we describe in the detail the baselines we use.

\subsection{Baselines}
\label{sec:baselines}
\textbf{Mistral} (Out of the Box LLM): Beam search is used to generate the Top-$K$ predictions for Mistral Instruct v0.2\\
\textbf{BM25}~\cite{robertson2009probabilistic}: We index content using Elastic Search v 8.11.2. Leaf node acts as a unique document identifier for a book\\
\textbf{DPR}: We use NV-Embed-v2\cite{lee2024nv} out of the box with text lengths set to 512 for passages and 256 for queries. NV-Embed-v2 is the leading text-only embedding model with open weights on the MTEB leaderboard\footnote{As per Jan 23'2025 MTEB Leaderboard Snapshot.}. The model is based on Mistral-7B-v0.1\footnote{https://huggingface.co/mistralai/Mistral-7B-v0.1}, making it a suitable baseline for {\sys} as both are finetuned over Mistral.\\
% \textbf{DPR}: We use NV-Embed-v2\cite{lee2024nv} out of the box with text lengths set to 512 for passages and 256 for queries. NV-Embed-v2 is the leading text-only embedding model with open weights on the MTEB leaderboard\footnote{As per Jan 23'2025 MTEB Leaderboard Snapshot.}. The model is based on Mistral-7B-v0.1\footnote{https://huggingface.co/mistralai/Mistral-7B-v0.1}, making it a suitable baseline for {\sys} as both are finetuned over Mistral. Unlike conventional methods for dense embeddings, such as mean pooling or <eos> token representation, NV-Embed uses a latent attention layer to obtain pooled embeddings, enhancing retrieval performance.\\
\textbf{DSI~\citep{Tay2022TransformerMA}} infuses the entire knowledge of a corpus in the parameters of an LLM and directly generates a document identifier. We fine-tune DSI using the train split for each book.

\subsection{Fine-tuning {\sys} and DSI}
\label{sec:finetune}
We fine-tune Mistral Instruct v0.2\footnote{https://huggingface.co/mistralai/Mistral-7B-Instruct-v0.2}~\cite{Jiang2023Mistral7} for a maximum of 200 epochs using a LoRA adapter~\cite{Hu2021LoRALA} with a $r=16$ and $\alpha=32$. We early stop with a patience of 20 epochs, by computing {\devmetric} on the dev set. We set the maximum input length to 14k tokens for {\sys} and $512$ tokens for DSI and a maximum output length of 64 tokens. See Table~\ref{table:train_example_small} for a sample training input and output for {\sys}.

% Table~\ref{table:train_prompt} lists the prompts used for fine-tuning {\sys} and DSI.

% \begin{table}[h]
% \caption{LLM Prompts for fine-tuning. {\sys} uses title, toc and query as prompt variables. DSI uses query as the prompt variable}
% %\small
% \begin{tabular}{l}
% {\sys}\\
% \toprule
% \#\#\# Book: \{title\}\\
% Pick the best section from the table of contents below \\ 
% which can answer the user query below. \\
% Only generate the section name, do not \\
% generate any explanation!\\
% \\
% \#\#\#\ Table of Contents\\
% \{toc\}\\
% \\
% \#\#\# Query: \{query\}\\
% \#\#\# Section: leaf\\
% \\

% \textbf{DSI}\\
% \midrule
% \#\#\# Query: \{query\}\\
% \#\#\# Section:
% \end{tabular}
% \label{table:train_prompt}
% \end{table}


% \subsection{Evaluation Metrics}
% We use the BeIR framework ~\cite{thakur2021beir} to evaluate our models and report Recall@1, Recall@3 and Normalised Discounted Cumulative Gain (NDCG@3) metrics on the test set for each book, after selecting the best model.